\subsection{Separation through degree \texorpdfstring{$n$}{n} in two alphabets}

A Hilbert-series identity alone need not eliminate cancelling homology.
Here multiplication already gives a free cover. Equality of dimensions
then forces its kernel to begin in a higher degree.

\begin{lemma}[Truncated freeness from a Hilbert identity]
\label{lem:truncated-freeness}
Let $S$ be a connected positively graded commutative or
supercommutative algebra over a field, and let $M$ be a nonnegatively
graded $S$-module. Assume their homogeneous pieces are finite
dimensional, and put $R=M/S_+M$. If
\[
 H_M(t)=H_S(t)H_R(t)\pmod{t^{N+1}},
\]
then every homogeneous section $R\to M$ induces an isomorphism
\[
 (S\otimes R)_{\le N}\longrightarrow M_{\le N}.
\]
Moreover, $\Tor_j^S(M,K)_d=0$ for $j>0$ and $d\le N$.
For superalgebras use ordinary dimensions and parity-preserving
sections. Additional positive multigradings may be retained.
\end{lemma}

\begin{proof}
The multiplication map is surjective by degree induction: after
subtracting a lift of the quotient class, an element belongs to
$S_+M$ and is a sum of positive-degree coefficients times
elements of smaller degree. The Hilbert identity makes the source
and target have equal dimensions through degree $N$. Its kernel
$L$ therefore has no homogeneous piece in those degrees.
Choose a graded free resolution of $L$ using generators only
in degrees greater than $N$, and splice it with
$0\to L\to S\otimes R\to M\to0$. Every positive term of the resulting
free resolution of $M$ is zero in degrees at most $N$.
Tensoring with $K$ proves the Tor assertion. All arguments respect
the additional gradings.
\end{proof}

\begin{theorem}[Bivariate separation through degree $n$]
\label{thm:bivariate-separation}
Multiplication induces a canonical $\GL_2\times S_n$-equivariant
isomorphism
\[
 \bigl(P_{n,2}^{S_n}\otimes\Harm_{n,2}\bigr)_{\le n}
          \xrightarrow{\ \sim\ }(P_{n,2})_{\le n}.
\]
Every positive moment Koszul homology group vanishes in internal
degrees at most $n$. Consequently the full bivariate $S_n$ character
is given, in bidegree $(a,b)$ whenever $a+b\le n$, by
\begin{equation}\label{eq:bivariate-character}
 \operatorname{tr}\bigl(\sigma\mid (R_{n,2})_{a,b}\bigr)
 =[q^at^b]\!
   \prod_{\substack{i,j\ge0\\i+j>0}}(1-q^it^j)
   \prod_{m\ge1}\bigl((1-q^m)(1-t^m)\bigr)^{-X_m(\sigma)}.
\end{equation}
Here $X_m(\sigma)$ counts the $m$-cycles of $\sigma$.
Equivalently, every row partition of length at most two has zero
positive moment homology through degree $n$, in every row dimension.
\end{theorem}

\begin{proof}
For every fixed bidegree $(a,b)$, with $D=a+b$, Wang's dimension
theorem \cite[Theorem~5.1]{Wang2026Polynomiality} gives one polynomial in $n$
equal to $\dim(R_{n,2})_{a,b}$ throughout $n\ge D$.
The Euler dimension of the moment complex is also polynomial in
$n$. Corollary~\ref{cor:universal-augmented-separation} gives eventual
acyclicity, for example for $n\ge2D-1$.
The actual and Euler dimension polynomials therefore agree at all
sufficiently large integers and are identical. Degrees zero,
one, and two follow from the independent linear and quadratic
moments. Hence the quotient has the Euler Hilbert series in every
total degree at most $n$.

Let
\[
 F_n=\Sym\left(\bigoplus_{j=1}^n\Sym^jK^2\right)
\]
act on $P_{n,2}$ by the formal moment evaluation. Polarized moments of
degrees at most $n$ generate the invariant algebra, so
$R_{n,2}=P_{n,2}/(F_n)_+P_{n,2}$.
The dimension
identity just proved is
$H_{P_{n,2}}=H_{F_n}H_{R_{n,2}}$ through degree $n$.
Lemma~\ref{lem:truncated-freeness} gives freeness and positive
Koszul acyclicity in this range.

The fundamental trace presentation has its first relation in
degree $n+1$: each relation uses $n+1$ positive inputs
\cite{Domokos2017}. Thus $F_n$ maps isomorphically
onto the actual invariant algebra in degrees at most $n$.
Choose the canonical harmonic section of $R_{n,2}$. The free
cover becomes exactly the multiplication map in the statement.

Taking equivariant characters of the moment complex gives
\eqref{eq:bivariate-character}. Each $m$-cycle contributes
$(1-q^m)^{-1}(1-t^m)^{-1}$ to the polynomial trace, and each
invariant moment contributes its Euler factor. Finally, the
complex is a polynomial functor of the row space. In
characteristic zero its homology decomposes into Schur functors;
every partition of length at most two is detected on $K^2$.
\end{proof}



